\chapter{Architektur und Konvergenzmechanismen von K3s und ArgoCD}
\label{chap:architektur-k3s-argocd}

Das untersuchte Zielsystem besteht aus der leichtgewichtigen Kubernetes-Distribution K3s als Laufzeitbasis und ArgoCD als deklarativer GitOps-Control-Plane. Um Latenzspitzen, Zustandsübergänge und Speicherprofile im Experiment kausal einordnen zu können, werden im Folgenden die Komponentenstruktur sowie die internen Polling-, Timeout- und Begrenzungsmechanismen beider Plattformen analysiert.

\section{Leichtgewichtige Container-Orchestrierung mit K3s}
\label{sec:leichtgewichtige-orchestrierung-k3s}

Standard-Kubernetes verteilt Control-Plane-Komponenten wie \texttt{kube-apiserver}, \texttt{etcd}, \texttt{kube-scheduler} und \texttt{kube-controller-manager} auf getrennte Prozesse und Container.\autocite{burns_borg_2016} Auf kompakten Edge-Knoten oder kleinen Cloud-Instanzen (etwa \gl{aws} \texttt{c7i-flex.large} mit 2~\gl{vcpu}s und 4~GiB \gl{ram}) belegt diese Basisschicht bereits einen spürbaren Teil des verfügbaren Hauptspeichers, noch bevor GitOps-Operatoren oder Monitoring-Dienste gestartet werden.\autocite{koziolek_k3s_2023}

Die \gl{cncf}-zertifizierte Distribution \textit{K3s} konsolidiert diesen Aufbau, indem sie sämtliche Control-Plane- und Worker-Komponenten samt \texttt{containerd} und dem \gl{cni}-Plugin \texttt{Flannel} in einer einzigen Binärdatei bündelt und veraltete Cloud-Provider-Treiber entfernt.\autocite{k3s_architecture_2026} Für Single-Node-Cluster ersetzt K3s zudem den speicherintensiven Key-Value-Store \texttt{etcd} standardmäßig durch eine eingebettete SQLite-Datenbank über die Abstraktionsschicht \textit{Kine}, die \texttt{etcd}-Aufrufe transparent in relationale Transaktionen übersetzt.\autocite{k3s_architecture_2026}

Empirische Vergleiche zeigen dabei, dass die Vorteile nicht pauschal in jeder Dimension gleich groß sind: Während Böhm und Wirtz vor allem deutliche Einsparungen beim Massenspeicher und schnellere Startzeiten feststellen, aber nur moderate \gl{ram}-Unterschiede zum regulären Kubernetes messen,\autocite{boehm_k3s_2021} dokumentieren Koziolek und Eskandani im Vergleich mehrerer Lightweight-Distributionen eine kompakte Leerlauf-Speicherauslastung von K3s.\autocite{koziolek_k3s_2023} Für kontrollierte Experimente lässt sich der Eigenbedarf weiter senken, indem vorinstallierte Zusatzdienste wie \textit{Traefik} und der \textit{Metrics-Server} bei der Installation deaktiviert werden (\texttt{--disable traefik --disable metrics-server}). Die volle \gl{api}- und \gl{crd}-Kompatibilität bleibt erhalten, während nahezu das gesamte Knotenbudget dem \gl{sut} und der Metrikerfassung zur Verfügung steht.

\section{Komponentenarchitektur und Zustandsmaschine von ArgoCD}
\label{sec:komponentenarchitektur-argocd}

ArgoCD trennt seine Control Plane in drei zentrale Microservices und einen internen Cache.\autocite{argocd_architecture_2026} Der \textit{\gl{api}-/Web-Server} (\texttt{argocd-\allowbreak{}server}) bietet \gl{grpc}-, \gl{rest}- und Web-\gl{ui}-Schnittstellen für Administration, Authentifizierung und manuelle Sync-Befehle, übernimmt aber keine periodischen Hintergrundabgleiche. Für die externe Git-Kommunikation ist der zustandsfreie \textit{Repository-Server} (\texttt{argocd-\allowbreak{}repo-\allowbreak{}server}) zuständig: Er klont konfigurierte Repositories, ruft Änderungen ab und rendert deklarative Vorlagen wie Kustomize-Overlays oder Helm-Charts in konkrete Kubernetes-Manifeste.\autocite{argocd_architecture_2026,kubernetes_kustomize_2026} Das eigentliche Regelungsherz bildet der Go-basierte \textit{Application Controller} (\texttt{argocd-application-controller}). Er beobachtet alle \texttt{Application}-Ressourcen, bezieht über den \texttt{repo-server} den Soll-Zustand, vergleicht diesen mithilfe des internen \texttt{redis}-Caches mit dem Live-Zustand im Cluster und führt bei Abweichungen Korrekturoperationen gegen den \texttt{kube-apiserver} aus.\autocite{argocd_architecture_2026}

Den Zustand jeder Anwendung modelliert der Controller über zwei orthogonale Dimensionen sowie eine Liste von \textit{Conditions} (Tabelle~\ref{tab:argocd-zustandsdimensionen}).

\begin{table}[H]
    \centering
    \small
    \begin{tabularx}{\textwidth}{>{\raggedright\arraybackslash}p{3.9cm}|>{\raggedright\arraybackslash}p{3.5cm}|>{\raggedright\arraybackslash}X}
    \textbf{Zustandsdimension} & \textbf{Mögliche Werte} & \textbf{Bedeutung im Reconcile-Zyklus} \\
    \specialrule{1.5pt}{0pt}{0pt}
    Sync-Status\newline (\texttt{\footnotesize .status.sync.status}) & \texttt{Synced}, \texttt{OutOfSync}, \texttt{Unknown} & Zeigt an, ob der Live-Zustand im Cluster mit den Git-Manifesten übereinstimmt. Scheitert der Git-Abruf, wechselt der Status ggf.\ in \texttt{Unknown}. \\
    \hline
    Health-Status\newline (\texttt{\footnotesize .status.health.status}) & \texttt{Healthy}, \texttt{Progressing}, \texttt{Degraded}, \texttt{Suspended}, \texttt{Missing}, \texttt{Unknown} & Bewertet die betriebliche Gesundheit der untergeordneten Kubernetes-Objekte (z.\,B.\ laufende Pod-Replicas eines \texttt{Deployment}). \\
    \hline
    Conditions\newline (\texttt{\footnotesize .status.conditions}) & u.\,a.\ \texttt{ComparisonError}, \texttt{SyncError}, \texttt{InvalidSpecError} & Protokolliert Fehler während des Abgleichs. Bei Netzwerkverlust zum Git-Server hinterlegt der Controller eine \texttt{ComparisonError}-Bedingung. \\
    \hline\hline
    \end{tabularx}
\caption{Zustandsdimensionen einer ArgoCD-\texttt{Application}-Ressource, eigene Darstellung nach Argo Project (2026).}
\label{tab:argocd-zustandsdimensionen}
\end{table}

Diese Trennung begründet ein zentrales Resilienzmerkmal: Selbst wenn eine Netzwerkstörung zum Git-Server eine \texttt{ComparisonError}-Bedingung auslöst, bleiben die bereits laufenden Produktiv-Workloads im Cluster \texttt{Healthy}. Solange der Controller stabil weiterarbeitet, bleibt der Fehler auf die Synchronisationsfähigkeit der Control Plane begrenzt.

\section{Reconcile-Zyklus, Timeout-Kaskaden und Ressourcen-Constraints}
\label{sec:reconcile-zyklus-timeouts-ressourcen}

Im Normalbetrieb pollt ArgoCD Git-Repositories periodisch, um Netzwerk- und \gl{cpu}-Last gering zu halten. Der Standardwert von \texttt{timeout.\allowbreak{}reconciliation} liegt bei \texttt{180s}; das eingesetzte Helm-Chart~\texttt{10.18.2} setzt \texttt{120s} zuzüglich bis zu \texttt{60s} Jitter (\texttt{timeout.\allowbreak{}reconciliation.\allowbreak{}jitter}).\autocite{argocd_reconcile_2026,argo_helm_2026} Zusätzlich cacht der \texttt{argocd-\allowbreak{}repo-\allowbreak{}server} gerenderte Manifeste in Redis: Bleibt die Commit-\gl{sha} bei einer \texttt{git ls-remote}-Abfrage unverändert, entfällt ein vollständiger \texttt{git fetch}- und Kustomize-Durchlauf.\autocite{argocd_reconcile_2026} Für ein zweiminütiges Störungsexperiment folgt daraus eine zentrale Designanforderung: Ohne aktiven Eingriff könnte das Chaos-Fenster vollständig zwischen zwei Polling-Zyklen liegen oder aus dem Cache bedient werden. Erst die Annotation \texttt{argocd.\allowbreak{}argoproj.io/\allowbreak{}refresh=hard} invalidiert den Manifest-Cache und erzwingt eine unmittelbare Git-Abfrage über den \texttt{repo-server}.\autocite{argocd_reconcile_2026}

Trifft eine solche Abfrage auf eine paketverlustbehaftete Strecke, greifen zwei technisch getrennte Mechanismen:

In der \textit{Reconciliation-Phase} ruft der \texttt{application-\allowbreak{}controller} per \gl{grpc} den \texttt{repo-server} auf, der seinerseits einen Git-Prozess über \gl{https} (Port~443) startet. Verzögern \gl{tcp}-Retransmissionen diesen Abruf, steigen sowohl die Dauer der Git-Anfrage (\texttt{argocd\_\allowbreak{}git\_\allowbreak{}request\_\allowbreak{}duration\_\allowbreak{}seconds}) als auch die gesamte Reconcile-Latenz (\texttt{argocd\_\allowbreak{}app\_\allowbreak{}reconcile}).\autocite{argocd_metrics_2026} Bricht der Git-Abruf per Timeout ab, setzt der Controller eine \texttt{ComparisonError}-Condition und überführt den \texttt{SyncStatus} in \texttt{Unknown}, ohne überhaupt eine Sync-Operation (\texttt{argocd\_\allowbreak{}app\_\allowbreak{}sync\_\allowbreak{}total}) gegen Kubernetes zu starten. Weitere Abgleichversuche werden in diesem Zustand über \gl{tcp}-/\gl{grpc}-Timeouts und die interne Workqueue- bzw.\ Polling-Planung gesteuert.

Die applikationsseitige \textit{Retry-Richtlinie} (\texttt{syncPolicy.retry.backoff}) greift dagegen erst dann, wenn bei bekanntem Soll-Zustand und erkanntem Drift (\texttt{OutOfSync}) eine aktive Auto-Sync-Operation (\texttt{selfHeal: true}) gegen den \texttt{kube-apiserver} selbst fehlschlägt.\autocite{argocd_autosync_2026} Bei \texttt{limit: 5}, Basiswartezeit $d_0 = 5\,\text{s}$, Faktor $f = 2$ und Deckel $d_{\max} = 180\,\text{s}$ berechnet sich die Wartezeit $d_n$ vor dem $n$-ten Wiederholungsversuch ($n \in \{1, \dots, 5\}$) zu:

\begin{equation}
d_n = \min(d_0 \cdot f^{n-1},\; d_{\max})
\label{eq:exponential-backoff}
\end{equation}

Die resultierende Staffelung ($5\,\text{s}, 10\,\text{s}, 20\,\text{s}, 40\,\text{s}, 80\,\text{s}$) verhindert bei anhaltenden Sync-Fehlern einen \textit{Retry Storm}.\autocite[Kap.~22, S.~259--272]{beyer_sre_2016} Bei reinen Git-Netzwerkstörungen bestimmen jedoch primär ablaufende \gl{tcp}-/\gl{grpc}-Timeouts sowie der Abstand bis zum nächsten Reconcile-Zyklus die Erholungszeit $\text{\gl{mttr}}_{\text{post-fault}}$.

Gleichzeitig halten wartende Reconcile-Worker im Controller sowie blockierte \gl{tls}-Kontexte und Kustomize-Kindprozesse im Repo-Server Arbeitsspeicher gebunden. Auf ressourcenbeschränkten Knoten werden beide Pods daher mit engen Speicherlimits versehen (\texttt{512Mi} für den Controller, \texttt{384Mi} für den Repo-Server). Überschreitet ein Container sein Limit, beendet das Linux-Cgroups-Subsystem den Prozess sofort mit \texttt{OOMKilled} (Exit-Code~137).\autocite[Kap.~5]{burns_kubernetes_2022} Resilienz bemisst sich somit nicht nur daran, ob ArgoCD nach einer Störung überhaupt konvergiert, sondern am Zusammenspiel aus Latenzanstieg, $\text{\gl{mttr}}_{\text{post-fault}}$ und dem Sicherheitsabstand zur \gl{oom}-Grenze bei voller Swap-Reinheit.
